\subsection*{The outer pressure datum as a positive measure}
\label{subsec:outer-pressure-measure}
The pressure datum of Lemma A.5 in \cite{openai2026} is defined using
the reference exterior schedule with its pressure-neutral angular
corrections omitted. The proof of that lemma states that each stage
has the form
\begin{equation}\label{eq:measure-profile}
 E_{\mathrm{id,sched}}(y,\eta)
   =c(y)(1+\eta^2)^{-\vartheta(y)},\qquad
 c(y)>0,\quad0\le\vartheta(y)\le1,
\end{equation}
where $c$ and $\vartheta$ are independent of $\eta$. Let $\mu$ be the
pushforward of $c(y)^2dy$ under $y\mapsto\vartheta(y)$, and put
$\mathfrak p=\frac12\mu([0,1])=-\Pi_0(0)>0$.

\begin{proposition}[Exact positive-measure pressure trace]
\label{prop:exact-pressure-measure}
The actual axis datum is exactly
\begin{equation}\label{eq:exact-pressure-measure}
 \Pi_0(\eta)=-\frac12\int_{[0,1]}
 (1+\eta^2)^{-2s}\,d\mu(s),\qquad
 \mu(\{1\})\ge5P_*^2.
\end{equation}
For $|\eta|\le1$,
\begin{equation}\label{eq:pressure-mass-bounds}
 -\mathfrak p\le\Pi_0(\eta)\le
 -\frac52P_*^2(1+\eta^2)^{-2},\qquad
 |\Pi_0'|\le2\mathfrak p,\quad
 |\Pi_0''|\le8\mathfrak p.
\end{equation}
Its first derivative has the strict sign of $\eta\ne0$.
The pressure extends holomorphically to
$|\operatorname{Re}z|\le11/10$,
$|\operatorname{Im}z|\le1/10$, with
\begin{equation}\label{eq:pressure-complex-bound}
 |\Pi_0(z)|\le(100/99)^2\mathfrak p,\qquad
 |\Pi_0^{(m)}(\eta)|\le(100/99)^2\mathfrak p\,m!20^m.
\end{equation}
\end{proposition}
\begin{proof}
The inner interval $y\le0$ has
$c(y)=P_*e^{y/10}$ and $\vartheta=1$; its contribution to
$\mu$ is $5P_*^2\delta_1$. The remaining stages contribute
a finite positive measure. For $w_s=(1+\eta^2)^{-2s}$,
\[
 \frac{w_s'}{w_s}=-\frac{4s\eta}{1+\eta^2},\qquad
 \frac{w_s''}{w_s}
 =-\frac{4s(1-\eta^2)}{(1+\eta^2)^2}
   +\frac{16s^2\eta^2}{(1+\eta^2)^2}.
\]
For $|\eta|\le1$, $|w_s'|\le2$,
$|w_s''|\le8$. Positivity of $\mu$ proves the real estimates;
the inner atom gives strictness of the derivative sign.
For complex $z$ in the stated rectangle,
$\operatorname{Re}(1+z^2)\ge99/100$; the principal complex
power is analytic and has modulus at most $(100/99)^2$.
Integrating over $\mu$ gives the complex bound.
Cauchy's estimate on radius-$1/20$ disks proves the last one.
\end{proof}

\subsection*{Quantitative choice of the B.2 parameters}
The full schedule is not numerically frozen, but the preceding
positivity permits explicit admissibility parameters depending
only on its pressure mass $\mathfrak p$ and the axis offset $j_0$.

\begin{proposition}[A constructive real-interval B.2 bound]
\label{prop:constructive-B2}
Assume the source Lemma A.5, $P_*\ge1$,
$0<h\le1/100$, $0<j_0\le1/20$. Set
\[
 U_*=4\eta+j_0,\quad D=\tfrac12-h,\quad
 A=\tfrac12+h,\quad d=1-\eta^2,\quad
 H_*=D\eta+dU_*,
\]
\[
 Z_*=-A(1-2\eta U_*)U_*-4H_*-d\Pi_0'
       +4A\eta\Pi_0.
\]
Let
\begin{equation}\label{eq:constructive-B2-parameters}
 K=200(1+\mathfrak p),\quad
 \delta_*=\frac{j_0}{12},\quad
 \epsilon=\frac{j_0}{12K},\quad
 \sigma_*=\frac{\epsilon}{30}.
\end{equation}
The unique zero $\eta_0$ of $H_*$ lies in
$(-j_0/4,-j_0/5)$. Moreover on $[-1,1]$,
\begin{equation}\label{eq:B2-certified}
 |Z_*|\le\delta_*
 \quad\Longrightarrow\quad |H_*|>\frac25\epsilon
 \quad\Longrightarrow\quad
 \frac{H_*^2}{H_*^2+\sigma_*^2}\ge\frac{144}{145}
 >\frac{99}{100}.
\end{equation}
These strict inequalities persist on some larger real compact
interval $I\supset[-1,1]$.
\end{proposition}
\begin{proof}
The quotient $H_*/d=D\eta/d+4\eta+j_0$
is strictly increasing on $(-1,1)$, and
$H_*(-j_0/4)=-Dj_0/4<0$ while
$H_*(-j_0/5)=(j_0/5)(1-D-j_0^2/25)>0$.
Write $t=-\eta_0$, so $j_0/5<t<j_0/4$.
At the root, $U_*(\eta_0)=Dt/(1-t^2)<j_0/7$
and $1-2\eta_0U_*(\eta_0)<2$.
The reference inner-tail inequality in
\eqref{eq:pressure-mass-bounds}, together with
$(1+t^2)^{-2}>1/2$, gives
$4A\eta_0\Pi_0(\eta_0)\ge5AtP_*^2\ge Aj_0P_*^2$.
Also $-d\Pi_0'(\eta_0)\ge0$. Therefore
\begin{equation}\label{eq:Z-root-positive}
 Z_*(\eta_0)>Aj_0(P_*^2-2/7)>j_0/3.
\end{equation}
Differentiating the exact formula for $Z_*$ and applying
\eqref{eq:pressure-mass-bounds} gives
$|Z_*'|\le120+20\mathfrak p<K$ on $[-1,1]$.
Indeed $|U_*|\le81/20$, $|H_*'|\le83/5$ and
$A\le51/100$ bound the polynomial terms.
Thus $|\eta-\eta_0|\le\epsilon$ implies
$Z_*(\eta)>j_0/3-K\epsilon=j_0/4>\delta_*$.

Factoring at the zero,
\[
 H_*(\eta)=(\eta-\eta_0)Q(\eta),\quad
 Q=D+4-4(\eta^2+\eta\eta_0+\eta_0^2)
      -j_0(\eta+\eta_0).
\]
The quadratic $Q$ is concave; on $[-1,1]$ its endpoint
values exceed $2/5$ because $D\ge49/100$ and
$|\eta_0|<j_0/4\le1/80$. Hence
$|H_*|>(2/5)|\eta-\eta_0|$. The condition
$|Z_*|\le\delta_*$ excludes $|\eta-\eta_0|\le\epsilon$.
Finally $\sigma_*=\epsilon/30$ yields $\chi\ge144/145$.
Compactness and strictness allow $I$ to be enlarged slightly.
\end{proof}

\paragraph{Unresolved numerical and dynamical data.}
The positive measure $\mu$ and its mass $\mathfrak p$ remain
defined by the actual full outer schedule and its moment matching,
not by an arbitrary model. Numerical evaluation requires freezing
that schedule and certified quadrature. The remaining nonlinear
Appendix-B construction requires effective values of $\Lambda$,
$C$, and a verified contraction ball. The theorem does not solve
the separate similarity-order $n=1,2$ unforced corrections or
establish the physical Navier--Stokes reconstruction.
